\section{Definitions}\label{sec:prelims}

\subsection{Key encapsulation mechanisms}\label{sec:kems}

\subsubsection{Syntax}

\begin{definition}[Key Encapsulation Mechanism (KEM)]
    \label{def:kem}
    A key encapsulation mechanism is a triple of algorithms $\kem =
    \{\keygen, \enc, \dec\}$ with public keyspace $\mathcal{PK}$, private
    keyspace $\mathcal{SK}$, ciphertext space $\mathcal{C}$, and shared
    keyspace $\mathcal{K}$.
    %
    The triple of algorithms is defined as:
    \begin{itemize}
    \item $\kem.\keygen(\, )\rsample (\sk, \pk)$ Randomized algorithm that
        outputs a secret (private) key $\sk \in \mathcal{SK}$, and a public
        key $\pk \in \mathcal{PK}$.
    \item $\kem.\enc(\pk)\rsample (k, c)$ Randomized algorithm that, given a
        public key $\pk \in \mathcal{PK}$, outputs a shared key $k \in
        \mathcal{K}$, and a ciphertext $c \in \mathcal{C}$.
    \item $\kem.\dec(c, \sk)\foutput y \in \{k, \bot\}$ Deterministic
        algorithm that, given a secret key, $\sk \in \mathcal{SK}$ and a
        ciphertext $c \in \mathcal{C}$, returns the shared key $k \in
        \mathcal{K}$.
        %
        In case of rejection, this algorithm returns $\bot$.
    \end{itemize}
\end{definition}


\subsubsection{Correctness}\label{sec:kem-correctness}
%
The correctness  of a KEM imposes  that, except with small  probability drawn
over the  random coin space of  $\kem.\keygen$ and $\kem.\enc$, we  have that
$\kem.\dec$ correctly recovers the shared key produced by $\kem.\enc$.
%
Formally, we say that a KEM is $\delta$-correct if:
\[
    \Pr \left[\,  k \ne k' \, \left | \substack{(\sk, \pk) \sample \keygen(\,) \\ (k, c) \sample \enc(\pk) \\ k' \gets \dec(c, \sk) } \right. \right] \le \delta .
\]


\subsubsection{Security}
%
The IND-CCA security game for KEMs is denoted as $\cca^b_{\kem,
    \adversary{A}}$ and is shown in Figure~\ref{fig:cca-game}.
%
\begin{gamefigure}{IND-CCA security game for KEMs.}{fig:cca-game}
    \begin{games}{2}
        \begin{gameenv}{\textbf{Game} $\cca^b_{\kem, \adversary{A}}$}
            \State $(\sk, \pk) \sample \kem.\keygen(\,)$
            \State $(k_0, c^*) \sample \kem.\enc(\pk)$
            \State $k_1 \sample \mathcal{K}$
            \State $b' \sample \adversary{A}^{\mathsf{Dec(\cdot)}}(\pk, c^*, k_b)$
            \State \Return $b'$
        \end{gameenv}
        \gamebreak
        \begin{gameenv}{\oracledec}
            \State $\ifthen{c \compare c^*}$
                \State $\quad $\Return $\bot$
                \State $k \gets \kem.\dec(c, \sk)$
                \State \Return $k$
            \end{gameenv}
        \end{games}
    \end{gamefigure}


    \begin{definition}[IND-CCA advantage for KEMs]
        The advantage of $\adversary{A}$ in breaking the IND-CCA security of
        $\kem$ is defined as
        %
        \[
            \advantage{\cca, \adversary{A}}{\kem} = \abs{\prone{\cca^0_{\kem, \adversary{A}}} - \prone{\cca^1_{\kem, \adversary{A}}}}.
        \]
        %
        A $\kem$ is IND-CCA secure for if for any PPT adversary
        $\adversary{A}$, $\advantage{\cca, \adversary{A}}{\kem}$ is
        negligible.
    \end{definition}

    \subsection{Pseudorandom function}\label{sec:prf}
    %
    A pseudorandom function (PRF) is a keyed deterministic function that cannot be distinguished from a truly random function.

    \begin{definition}[Pseudorandom Function (PRF)]
        \label{def:prf}
        For a finite keyspace $\mathcal{K}$, input space $\mathcal{X}$,
        finite output space $\mathcal{Y}$, and an efficiently computable
        function $f: \mathcal{K} \times \mathcal{X} \foutput \mathcal{Y}$, we
        define the PRF advantage of adversary $\adversary{A}$ based on the
        experiments in Figure \ref{fig:prf-game} as
        %
        \[
            \advantage{\prf, \adversary{A}}{f} = \abs{\prone{\prf^0_{f, \adversary{A}}} - \prone{\prf^1_{f, \adversary{A}}}}.
        \]
        %
    \end{definition}

    \begin{gamefigure}{PRF security games.}{fig:prf-game}
        \begin{games}{2}
            \begin{gameenv}{\textbf{Game} $\prf^b_{f, \adversary{A}}$}
                \State $\pset{T}[\cdot] \assign \bot$
                \State $k \sample \mathcal{K}$
                \State $b' \sample \adversary{A}^{\Eval(\cdot)}$
                \State \Return $b'$
            \end{gameenv}
            \gamebreak
            \begin{gameenv}{\oracle $\Eval_{f}(x)$}
                \State $\ifthen{x \in \pset{T}}$
                    \State $\quad $\Return $\pset{T}[x]$
                    \State $y_0 \assign f(k, x)$; $y_1 \sample \mathcal{Y}$
                    \State $\pset{T}[x] \assign y_b$
                    \State \Return $y_b$
                \end{gameenv}
            \end{games}
        \end{gamefigure}
        %
        \subsection{Nominal Group}
        %
        In this section, we recall the definition of a nominal group, as of \cite{EC:ABHKLR21}.
        %
        \begin{definition}[Nominal Group \cite{EC:ABHKLR21}]
            A nominal group $\nominalgroup =\nominalgroupexpanded$ consists
            of an efficiently recognizable finite set of elements
            $\mathcal{G}$ (also called ``group elements''), a base element
            $\ngen \in \mathcal{G}$, a prime $p$, a finite set of honest
            exponents $\honestexponent \subset \mathbb{Z}$, a finite set of
            exponents $\exponent \subset \mathbb{Z} / p\mathbb{Z}$, and an
            efficiently computable exponentiation function $\nexp :
            \mathcal{G} \times \mathbb{Z} \foutput \mathcal{G}$, where we
            write $X^y$ for $\nexp(X, y)$.
            %
            The exponentiation function is required to have the following
            properties:
            %
            \begin{enumerate}
            \item $(X^y)^z = X^{yz} $ for all $ X \in \mathcal{G}, y, z \in
                \mathbb{Z}$.
            \item The function $\phi$ defined by $\phi(x) = g^x$ is a
                bijection from $\exponent$ to $\{\ngen^x | x \in [1, p -
                1]\}$.
            \end{enumerate}
        \end{definition}
        %
        As done in \cite{EC:ABHKLR21} for a nominal group $\nominalgroup =
        \nominalgroupexpanded$, we define $D_H$ to be the uniform
        distribution of honestly generated exponents $\honestexponent$ and
        $D_U$ to be the uniform distribution on $\exponent$.
        %
        We also recall the definition for the statistical distance between
        these distributions: $\Delta_{\nominalgroup} = \Delta [ D_H, D_U ] =
        \frac{1}{2} \underset{x \in \mathbb{Z}}{\sum} | \underset{D_U}{Pr}(x)
        - \underset{D_H}{Pr}(x) | $ \cite{EC:ABHKLR21}.
        %
        \subsection{Strong Diffie-Hellman Problem}
        \label{sec:strong-dh}

        \begin{definition}[Strong Diffie-Hellman (\sdh) Problem]
            %
            We define the advantage function of an adversary $\adversary{A}$
            against the Strong Diffie-Hellman problem over the nominal group
            $\nominalgroup$ as
            %
            \[
                \advantage{\nominalgroup, \adversary{A}}{\sdh} = \Pr \left [Z = g^{xy} \left | x,y \assign \varepsilon_u; Z \assign \adversary{A}^{DH(\cdot, \cdot)}(g^x, g^y) \right. \right].
            \]
            %
            where DH is a decision oracle that on input $(Y, Z)$ with $Y, Z
            \in \mathcal{G}$, returns $1$ iff $Y^x = Z$ and $0$ otherwise.
        \end{definition}
